%! $n$th Roots \& Rational Exponents — Unit 6, Lesson 6.1 — Guided Notes
\documentclass[10pt]{article}
\usepackage{algebra2-article}
\usepackage{algebra2-boxes}
\newcommand{\TallMath}[1]{$\displaystyle #1\rule[-1.4em]{0pt}{3.2em}$}
\newcommand{\lgrid}[4]{%
  \draw[step=1, linegray!40, very thin] (#1,#3) grid (#2,#4);
  \draw[->, semithick, charcoal] (#1,0)--(#2,0) node[below right, font=\scriptsize]{$x$};
  \draw[->, semithick, charcoal] (0,#3)--(0,#4) node[left, font=\scriptsize]{$y$};}
% Cube root: pgfmath has no cbrt, so plot the two branches separately.
\newcommand{\cbrtpos}[2]{\draw[very thick, forest, domain=#1:#2, samples=70, smooth]
  plot (\x,{pow(\x,1/3)});}
\newcommand{\cbrtneg}[2]{\draw[very thick, forest, domain=#1:#2, samples=70, smooth]
  plot (\x,{-pow(-1*(\x),1/3)});}

\begin{document}

\pageheader{Unit 6, Lesson 6.1}{Guided Notes}
\namedateperiod

\vspace{0.08in}

\begin{objectivebox}
\textbf{By the end of this lesson, I will be able to\dots}
\begin{itemize}\setlength{\itemsep}{2pt}
  \item find \emph{$n$th roots}, name the \emph{principal} root, and explain --- with exponents --- why an even root of a negative number does not exist while an odd one does;
  \item convert back and forth between \emph{radical form} and \emph{rational exponent form}, using $a^{m/n}=\sqrt[n]{a^m}=\left(\sqrt[n]{a}\right)^m$, and evaluate such expressions without a calculator;
  \item apply the exponent properties I already know to rational exponents --- including combining radicals with \emph{different indices}, which no radical rule can do.
\end{itemize}
\end{objectivebox}

\vspace{0.05in}

\begin{vocabbox}
\small
Fill in each term as we build it together.
\par\vspace{2pt}   % \par is required: \termblanklong's \noindent is a no-op mid-paragraph
\termblanklong{$n$th root}
\termblanklong{Principal $n$th root}
\termblanklong{Rational exponent}
\termblanklong{The $a^{m/n}$ rule}
\termblanklong{Negative rational exponent}
\end{vocabbox}

\begin{hookbox}
Here is a problem you cannot do --- yet:
\[
  \sqrt{x}\cdot\sqrt[3]{x} = \underline{\hspace{2.5cm}}
\]
Every radical rule you have ever seen requires the two indices to \textbf{match}, and these do not.
\begin{itemize}\setlength{\itemsep}{1pt}
  \item So try the other direction. What number, raised to the power $2$, gives $a$? \writeline
  \item Now suppose $a^{1/2}$ meant something. Using ``power to a power,'' $\left(a^{1/2}\right)^2 = a^{\,\blank{1.2cm}} = a$. So $a^{1/2}$ has \textbf{no choice} --- it must equal \blank{2.2cm}.
\end{itemize}
That is today's whole idea: \textbf{a radical is an exponent in disguise.} Once you can write
$\sqrt[n]{a}$ as $a^{1/n}$, every exponent rule from the Warm-Up applies to radicals --- and the
problem above becomes one line.
\end{hookbox}

\begin{notesbox}{1. $n$th Roots and the \emph{Principal} Root}
\textbf{Definition.} $\sqrt[n]{a} = b$ means $b$ raised to the \blank{1.2cm} power equals $a$.
\; The \textbf{index} $n$
says which root; the number under the radical is the \textbf{radicand}.

\vspace{3pt}
But ``which root'' hides a subtlety. Since $4^2=16$ \emph{and} $(-4)^2=16$, the number $16$ has
\emph{two} real square roots. The radical symbol is defined to give only the
\textbf{principal} (the \emph{non-negative}) one:
\[
  \sqrt{16} = \blank{1.2cm}, \qquad -\sqrt{16} = \blank{1.2cm},
  \qquad \pm\sqrt{16} = \blank{1.6cm}.
\]
If you want both roots, you must \emph{write} the $\pm$ yourself --- the symbol will not give it to you.

\vspace{3pt}
\textbf{How many real $n$th roots are there?} Fill in the table from the Warm-Up.
\begin{center}
\renewcommand{\arraystretch}{1.35}
\begin{tabularx}{\linewidth}{@{}l l X X@{}}
\hline
\textbf{Index $n$} & \textbf{Radicand $a$} & \textbf{How many real $n$th roots?} & \textbf{Example} \\
\hline
\textbf{Even} & $a>0$ & \blank{2.0cm} (one positive, one negative) & $\sqrt[4]{81}=\blank{0.8cm}$, and $-3$ also works \\
\textbf{Even} & $a=0$ & \blank{2.0cm} & $\sqrt{0}=0$ \\
\textbf{Even} & $a<0$ & \blank{2.0cm} & $\sqrt{-16}$ is \blank{2.4cm} \\
\textbf{Odd}  & any $a$ & \blank{2.0cm}, with the same sign as $a$ & $\sqrt[3]{-27}=\blank{0.8cm}$ \\
\hline
\end{tabularx}
\end{center}
\textbf{Why?} Warm-Up item~1 already answered it. An \textbf{even} power of any real number is
\blank{3.0cm}, so nothing real can be raised to an even power and come out negative. An
\textbf{odd} power \blank{3.0cm} the sign, so a negative radicand is no problem at all. \emph{This
is the algebraic reason behind Lesson 6.0's domain rule.}

\vspace{3pt}
\textbf{Perfect $n$th powers to know on sight} (you will read these \emph{backwards} all period):
\begin{center}
\renewcommand{\arraystretch}{1.3}
\begin{tabularx}{\linewidth}{@{}l X@{}}
Squares & $4,\ 9,\ 16,\ 25,\ 36,\ 49,\ 64,\ 81,\ 100,\ 121,\ 144$ \\
Cubes   & $8,\ 27,\ 64,\ 125,\ 216$ \quad and $-8,\ -27,\ -64,\dots$ \\
Fourth powers & $16,\ 81,\ 256$ \qquad Fifth powers: $32,\ 243$ \\
\end{tabularx}
\end{center}
\begin{itemize}[leftmargin=1.5em, itemsep=2pt]
  \item $\sqrt[5]{32}=\blank{1.0cm}$ \quad because $\blank{0.8cm}^5=32$.
  \item $\sqrt[4]{81}=\blank{1.0cm}$ \quad $\sqrt[3]{-64}=\blank{1.0cm}$ \quad
        $\sqrt[3]{\dfrac{8}{125}}=\blank{1.2cm}$ \quad (take the root of the top \emph{and} the bottom).
\end{itemize}
\end{notesbox}

\begin{notesbox}{2. The Trap: $\sqrt[n]{a^n}$ Is Not Always $a$}
Simplify each one \emph{carefully} --- work from the inside out.
\begin{center}
\renewcommand{\arraystretch}{1.5}
\begin{tabular}{@{}l l l@{}}
(a) \; $\sqrt{5^2}=\sqrt{\rule{0pt}{1.7ex}\blank{1.0cm}}=\blank{1.0cm}$ &
(b) \; $\sqrt{(-5)^2}=\sqrt{\rule{0pt}{1.7ex}\blank{1.0cm}}=\blank{1.0cm}$ &
(c) \; $\sqrt[3]{(-5)^3}=\sqrt[3]{\rule{0pt}{1.7ex}\blank{1.4cm}}=\blank{1.0cm}$ \\
\end{tabular}
\end{center}
In (b) the answer came back \emph{positive} even though we started with $-5$. The square destroyed
the sign, and the principal root refuses to put it back. So:
\begin{center}
\begin{tabular}{@{}l l@{}}
$n$ \textbf{even}: \; $\sqrt[n]{a^n}=\blank{1.6cm}$ &
\qquad $n$ \textbf{odd}: \; $\sqrt[n]{a^n}=\blank{1.6cm}$ \\
\end{tabular}
\end{center}
So $\sqrt{x^2}=\blank{1.4cm}$ (we do not know whether $x$ is positive!), while
$\sqrt[3]{x^3}=\blank{1.4cm}$.
\end{notesbox}

\begin{notesbox}{3. A Radical \emph{Is} an Exponent: $\sqrt[n]{a}=a^{1/n}$}
Nobody decided this --- the exponent rules \textbf{force} it. Whatever $a^{1/n}$ means, ``power to a
power'' has to keep working:
\[
  \left(a^{1/n}\right)^n = a^{\frac{1}{n}\cdot n} = a^{\blank{0.8cm}} = a .
\]
So $a^{1/n}$ is the number that, raised to the $n$th power, gives back $a$ --- and that is exactly
the definition of \blank{3.2cm}. Therefore:
\begin{center}
\begin{tcolorbox}[colback=goldbg, colframe=goldacc, boxrule=0.8pt, arc=1.5mm,
    left=3mm, right=3mm, top=1.5mm, bottom=1.5mm, width=0.62\linewidth]
\centering $\displaystyle \sqrt[n]{a} = a^{1/n}$ \qquad \emph{(the index becomes the denominator)}
\end{tcolorbox}
\end{center}
\begin{itemize}[leftmargin=1.5em, itemsep=3pt]
  \item $\sqrt{7}=7^{\,\blank{0.9cm}}$ \quad $\sqrt[3]{x}=x^{\,\blank{0.9cm}}$ \quad
        $\sqrt[5]{2y}=(2y)^{\,\blank{0.9cm}}$
  \item $11^{1/2}=\blank{1.6cm}$ \quad $m^{1/4}=\blank{1.6cm}$ \quad
        $64^{1/3}=\blank{1.0cm}$ (and now you can \emph{evaluate} it: \blank{0.8cm})
\end{itemize}
\emph{Watch the parentheses:} $2y^{1/5}$ and $(2y)^{1/5}$ are different --- in the first, only the
$y$ is under the root.
\end{notesbox}

\begin{notesbox}{4. The Full Rule: $a^{m/n}$ --- Denominator Roots, Numerator Powers}
What if the numerator is not $1$? Split the fraction: $\dfrac{m}{n}=\dfrac{1}{n}\cdot m$, so
$a^{m/n}=\left(a^{1/n}\right)^m$. Both orders are legal:
\begin{center}
\begin{tcolorbox}[colback=goldbg, colframe=goldacc, boxrule=0.8pt, arc=1.5mm,
    left=3mm, right=3mm, top=1.5mm, bottom=1.5mm, width=0.86\linewidth]
\centering $\displaystyle a^{m/n} = \sqrt[n]{a^m} = \left(\sqrt[n]{a}\right)^m$
\qquad \textbf{denominator} $=$ the \blank{1.8cm}, \quad \textbf{numerator} $=$ the \blank{1.8cm}
\end{tcolorbox}
\end{center}

\vspace{2pt}
\textbf{Converting (both directions).}
\begin{center}
\renewcommand{\arraystretch}{1.5}
\begin{tabularx}{\linewidth}{@{}X X | X X@{}}
\textbf{Radical form} & \textbf{Rational exponent} & \textbf{Rational exponent} & \textbf{Radical form} \\
\hline
$\sqrt[4]{x^3}$ & $x^{\,\blank{0.9cm}}$ & $y^{2/5}$ & \blank{2.0cm} \\
$\left(\sqrt[3]{x}\right)^7$ & $x^{\,\blank{0.9cm}}$ & $w^{5/2}$ & \blank{2.0cm} \\
$\sqrt{x^5}$ & $x^{\,\blank{0.9cm}}$ & $(3t)^{3/4}$ & \blank{2.0cm} \\
\end{tabularx}
\end{center}
\emph{The single most common error of this lesson is flipping the fraction.} $\sqrt[4]{x^3}$ is
$x^{3/4}$, \textbf{not} $x^{4/3}$ --- the \emph{index} is on the bottom, because it came from
$\frac{1}{n}$.

\vspace{4pt}
\textbf{Evaluating: take the root \emph{first}.} Both orders give the same answer, but one is
arithmetic you can do in your head.
\begin{center}
\renewcommand{\arraystretch}{1.6}
\begin{tabularx}{\linewidth}{@{}l X X@{}}
\hline
 & \textbf{Power first (hard)} & \textbf{Root first (easy)} \\
\hline
$8^{2/3}$ & $\sqrt[3]{8^2}=\sqrt[3]{\rule{0pt}{1.7ex}\blank{1.2cm}}=\blank{1.0cm}$ &
            $\left(\sqrt[3]{8}\right)^2=\left(\blank{0.8cm}\right)^2=\blank{1.0cm}$ \\
$16^{3/4}$ & $\sqrt[4]{16^3}=\sqrt[4]{\rule{0pt}{1.7ex}\blank{1.6cm}}=\blank{1.0cm}$ &
             $\left(\sqrt[4]{16}\right)^3=\left(\blank{0.8cm}\right)^3=\blank{1.0cm}$ \\
\hline
\end{tabularx}
\end{center}

\vspace{2pt}
\textbf{Negative rational exponents} mean exactly what they meant in the Warm-Up ---
\textbf{reciprocal}, never ``negative answer'':
\[
  a^{-m/n} = \frac{1}{a^{\,m/n}} .
\]
\begin{itemize}[leftmargin=1.5em, itemsep=3pt]
  \item $16^{-1/2} = \dfrac{1}{16^{\,1/2}} = \dfrac{1}{\blank{0.8cm}} = \blank{1.0cm}$ \quad
        (\emph{not} $-4$ --- the minus sign moves the expression, it does not change its sign).
  \item $27^{-2/3} = \dfrac{1}{\left(\sqrt[3]{27}\right)^{2}} = \dfrac{1}{\blank{0.8cm}}
        = \blank{1.0cm}$
\end{itemize}
\end{notesbox}

\begin{notesbox}{5. The Payoff: Every Exponent Property Now Works on Radicals}
The four rules from the Warm-Up never cared whether the exponent was a whole number. Rewrite, apply,
and (if asked) convert back.
\begin{center}
\renewcommand{\arraystretch}{1.4}
\begin{tabularx}{\linewidth}{@{}l X@{}}
\hline
Product & $a^{r}\cdot a^{s}=a^{\,\blank{1.4cm}}$ \\
Quotient & $\dfrac{a^{r}}{a^{s}}=a^{\,\blank{1.4cm}}$ \\
Power to a power & $\left(a^{r}\right)^{s}=a^{\,\blank{1.4cm}}$ \\
Negative & $a^{-r}=\blank{1.4cm}$ \\
\hline
\end{tabularx}
\end{center}
Now return to the Hook. Two \emph{different} indices, which no radical rule can touch:
\[
  \sqrt{x}\cdot\sqrt[3]{x} \;=\; x^{1/2}\cdot x^{1/3}
  \;=\; x^{\,\blank{1.4cm}} \;=\; \blank{2.4cm}\ \text{(back in radical form)}.
\]
\begin{itemize}[leftmargin=1.5em, itemsep=3pt]
  \item $\dfrac{x^{3/4}}{x^{1/2}} = x^{\,\blank{1.4cm}}$ \hfill
        $\left(x^{2/3}\right)^{6} = x^{\,\blank{1.4cm}}$
  \item $\sqrt[4]{x}\cdot\sqrt[4]{x^{2}} = x^{\,\blank{1.4cm}}$ \hfill
        $\left(8x^{6}\right)^{1/3} = \blank{2.2cm}$
\end{itemize}
\emph{This is why the notation exists.} Radicals with different indices are a dead end; the same
expressions as exponents are a one-line problem.
\end{notesbox}

\begin{notesbox}{6. Back to Lesson 6.0: the Denominator Decides the Domain}
These are the same two graphs you read yesterday --- only the notation changed.
\begin{center}
\begin{tikzpicture}[scale=0.34]
  % y = x^(1/2)
  \begin{scope}
    \lgrid{-2}{9}{-2}{4}
    \begin{scope}
      \clip (-2,-2) rectangle (9,4);
      \draw[very thick, forest, domain=0:9, samples=90, smooth] plot (\x,{sqrt(\x)});
    \end{scope}
    \fill[forest] (0,0) circle (5pt);
    \node[charcoal, font=\scriptsize] at (3.5,-3.4) {$y=x^{1/2}$};
  \end{scope}
  % y = x^(1/3)
  \begin{scope}[xshift=19cm]
    \lgrid{-8}{8}{-3}{3}
    \begin{scope}
      \clip (-8,-3) rectangle (8,3);
      \cbrtpos{0}{8}
      \cbrtneg{-8}{0}
    \end{scope}
    \node[charcoal, font=\scriptsize] at (0,-4.4) {$y=x^{1/3}$};
  \end{scope}
\end{tikzpicture}
\end{center}
\begin{itemize}[leftmargin=1.5em, itemsep=3pt]
  \item $y=x^{1/2}$ is another name for $y=\blank{1.6cm}$, so its domain is \blank{2.4cm}.
  \item $y=x^{1/3}$ is another name for $y=\blank{1.6cm}$, so its domain is \blank{2.6cm}.
  \item \textbf{The rule, in exponent language:} in $x^{m/n}$ (in lowest terms), if the
        \textbf{denominator} $n$ is \blank{1.6cm}, the domain is restricted to $x\ge 0$;
        if $n$ is \blank{1.6cm}, the domain is all reals. The \emph{numerator} never restricts
        anything.
  \item Test it: $x^{2/5}$ --- denominator \blank{0.8cm}, so the domain is \blank{2.4cm}. \quad
        $x^{3/4}$ --- denominator \blank{0.8cm}, so the domain is \blank{2.4cm}.
\end{itemize}
\emph{``Even index restricts the domain'' and ``even denominator restricts the domain'' are the same
sentence.}
\end{notesbox}

\begin{practicebox}
Work these with your class. Assume all variables represent positive real numbers.
\begin{enumerate}[leftmargin=1.7em, itemsep=5pt]
  \item Write in rational exponent form: \; $\sqrt[5]{x^{2}}=\blank{1.6cm}$ \quad
        $\left(\sqrt{y}\right)^{3}=\blank{1.6cm}$ \quad $\sqrt[3]{7t}=\blank{2.0cm}$
  \item Write in radical form: \; $x^{4/3}=\blank{2.0cm}$ \quad $c^{1/6}=\blank{2.0cm}$ \quad
        $(5n)^{3/2}=\blank{2.4cm}$
  \item Evaluate without a calculator (root first!): \; $125^{2/3}=\blank{1.4cm}$ \quad
        $81^{3/4}=\blank{1.4cm}$ \quad $32^{-2/5}=\blank{1.4cm}$
  \item Simplify using the exponent properties; leave the answer with a positive rational exponent:
        \; $x^{1/4}\cdot x^{1/2}=\blank{1.6cm}$ \quad $\dfrac{y^{5/6}}{y^{1/3}}=\blank{1.6cm}$ \quad
        $\left(w^{3/4}\right)^{2/3}=\blank{1.6cm}$
  \item \textbf{Explain.} A classmate writes $\sqrt[3]{x^{5}} = x^{3/5}$. What did they do wrong, and
        how would you remind them which number goes on the bottom? \writelines{2}
\end{enumerate}
\end{practicebox}

\end{document}
